\section{Funkcja nagrody}
\label{sec:reward-functions}

W obu wariantach zastosowano nagrodę terminalną. Kroki
niekończące rozdania zwracają wartość zero, natomiast po zakończeniu
epizodu nagroda wynika bezpośrednio z finansowego rozliczenia
rozdania. Dzięki temu optymalizowany cel odpowiada rzeczywistemu
wynikowi strategii, bez dodatkowego premiowania określonych akcji
lub układów kart.

Pierwszy wariant, Reward A, wykorzystuje zysk netto gracza wyrażony
w jednostkach stawki Ante:

\begin{equation}
    R_A =
    \begin{cases}
        0, & \text{dla kroku nieterminalnego}, \\
        P_{\mathrm{net}}, & \text{dla kroku terminalnego}.
    \end{cases}
\end{equation}

Drugi wariant, Reward B, stanowi liniowo przeskalowaną wersję
pierwszego wariantu:

\begin{equation}
    R_B = \frac{R_A}{6}.
\end{equation}

Wartość sześć odpowiada największej standardowej sumie stawek
Ante, Blind i Play w pojedynczym rozdaniu. Skalowanie nie jest
clippingiem, dlatego przy wysokich wypłatach zakładu Blind
nagroda nadal może przekraczać wartość jeden.

Oba warianty opisują ten sam cel ekonomiczny i zachowują
uporządkowanie wyników. Różnią się wyłącznie skalą sygnału
przekazywanego do algorytmu. Pozwala to zbadać wpływ skali nagrody
na przebieg procesu uczenia bez zmiany optymalizowanego celu.